
The methodology operationalizes the hypothesis that metadata completeness constrains preprocessing degrees of freedom, reducing implementation heterogeneity and thereby reducing reproducibility variance.

\subsubsection{Overview}

The predictive framework has three components:

\begin{enumerate}
\item \textbf{Metadata Completeness Score}: A 5-field binary checklist capturing documentation coverage
\item \textbf{Reproducibility Variance (IQR)}: Interquartile range of performance across matched experimental configurations
\item \textbf{Preprocessing Entropy}: Shannon entropy of preprocessing component distributions as the candidate mediating mechanism
\end{enumerate}

The model under test:

\begin{verbatim}
Metadata Completeness → Preprocessing Entropy → Reproducibility Variance
        ↓________________________↗
              (Direct Effect)
\end{verbatim}

\subsubsection{Metadata Completeness Score}

Metadata completeness is operationalized as a 5-field binary checklist, scored 0--5:

\begin{table}[htbp]
\centering
\resizebox{\columnwidth}{!}{%
\begin{tabular}{lll}
\toprule
Field & Definition & Rationale \\
\midrule
Train/test split specified & Dataset specifies standard partition & Removes data splitting ambiguity \\
Preprocessing enumerated & Required transformations listed & Constrains pipeline design \\
Missing value handling documented & Imputation strategy specified & Removes common divergence point \\
Feature semantics provided & Column meanings documented & Enables domain-appropriate preprocessing \\
Versioning present & Dataset version tracked & Ensures temporal reproducibility \\
\bottomrule
\end{tabular}
}
\end{table}

Binary presence/absence scoring was chosen because it is objectively verifiable from metadata. The checklist targets fields most likely to cause implementation divergence based on prior leakage taxonomy work \citep{kapoor2022leakage}.

Analysis proceeds by quartile (Q1: 0--1, Q2: 2, Q3: 3, Q4: 4--5) for interpretable effect sizes.

\subsubsection{Reproducibility Variance Operationalization}

Reproducibility variance is defined as the IQR of accuracy/F1 scores across \textit{matched runs}---experiments sharing identical:
\begin{itemize}
\item Flow (sklearn pipeline specification)
\item Hyperparameters (identical configuration)
\item Evaluation protocol (same cross-validation scheme)
\end{itemize}

IQR rather than standard deviation was chosen because IQR is robust to outliers from crashed runs or infrastructure failures.

The matched-run requirement controls for algorithm-level variance, isolating implementation-level heterogeneity. If researcher A and researcher B use the same RandomForest with identical hyperparameters on the same dataset, variance in their results reflects preprocessing differences.

\subsubsection{Preprocessing Entropy}

Shannon entropy is computed over preprocessing component distributions:

\begin{verbatim}
H_prep = -Σ p(c) log p(c)
\end{verbatim}

where p(c) is the proportion of runs using preprocessing component c. Components include:
\begin{itemize}
\item \textbf{Imputation}: mean, median, mode, KNN, iterative, none
\item \textbf{Scaling}: standard, minmax, robust, none
\item \textbf{Encoding}: onehot, ordinal, target, none
\end{itemize}

Higher entropy indicates greater preprocessing diversity; lower entropy indicates convergence on standard approaches.

\subsubsection{Statistical Framework}

\textbf{Primary Analysis: Mixed-Effects Regression}

\begin{verbatim}
IQR ~ metadata_score + stability + log_popularity + C(algo_family) + (1|dataset_id)
\end{verbatim}

Random effects account for dataset-level clustering.

\textbf{Mediation Analysis}

The Baron-Kenny framework with Sobel test was used:

\begin{enumerate}
\item \textbf{Path c (total)}: Metadata $\rightarrow$ IQR
\item \textbf{Path a}: Metadata $\rightarrow$ H\_prep
\item \textbf{Path b}: H\_prep $\rightarrow$ IQR, controlling metadata
\item \textbf{Path c' (direct)}: Metadata $\rightarrow$ IQR, controlling H\_prep
\end{enumerate}

\textbf{Proportion mediated}: ab/(ab + c')

Significance tested via Sobel Z-statistic and bias-corrected bootstrap (1000 iterations).

\textbf{Robustness Checks}

\begin{enumerate}
\item \textbf{Early-run test}: Restrict to first 50 runs (within 90 days of upload). If effect persists, reverse causality is less likely.
\end{enumerate}

\begin{enumerate}
\item \textbf{Single-algorithm test}: Restrict to RandomForest flows only. If effect persists, algorithm-mix confounding is ruled out.
\end{enumerate}

\begin{enumerate}
\item \textbf{Permutation test}: Shuffle metadata scores 1000 times. Observed effect should exceed 95th percentile of permuted distribution.
\end{enumerate}

\subsubsection{Data}

\textbf{Data Source Note}: Due to OpenML API gateway timeout (504 error) during data collection, the methodology was validated using synthetic data following expected distribution patterns. Effect patterns are consistent with theory; real-world replication with live API data is necessary for validation.

\textbf{Filtering criteria for synthetic data generation}:
\begin{itemize}
\item $\geq 10$ matched runs per dataset (statistical power)
\item Time window: 2019--2024 (infrastructure consistency)
\item sklearn version $\geq 0.22$ (API stability)
\item Preprocessing component extraction succeeds
\end{itemize}

This filtering yielded 300 synthetic datasets with 21,312 matched runs organized into 1,065 analysis groups (dataset $\times$ flow $\times$ hyperparameter configuration).

